% 第5章 5.4.2：最小值核 min(x,z)（[0,1]、Lebesgue 测度）下两种范数的谱权重（解析式）。
% (a) f=psi_1 与 g=psi_1+0.1 psi_8，psi_j=sqrt2 sin((j-1/2)pi x)：||f-g||_{L2}=0.1，||f-g||_{H_K}=0.1*7.5pi=2.356。
% (b) L^2 权重 1 与 RKHS 权重 1/lambda_j=((j-1/2)pi)^2（对数纵轴）。
\begin{tikzpicture}[font=\figurefont\footnotesize]
\pgfplotsset{every axis/.style={width=52mm,height=40mm,axis lines=left,
  x axis line style={draw=FigureStroke,line width=.6pt},y axis line style={draw=FigureStroke,line width=.6pt},
  tick style={draw=FigureStroke},tick label style={font=\figurefont\scriptsize,text=FigureInk},
  label style={font=\figurefont\scriptsize},
  ymajorgrids,grid style={draw=FigureGrid,line width=.4pt},
  title style={font=\figurefont\footnotesize},
  legend style={draw=none,fill=white,fill opacity=.85,text opacity=1,font=\figurefont\scriptsize,cells={anchor=west}}}}
\begin{axis}[name=A,xmin=0,xmax=1,ymin=-0.1,ymax=1.65,xtick={0,0.5,1},ytick={0,0.5,1,1.5},
  xlabel={$x$},xlabel style={at={(axis description cs:1,0)},anchor=north west,yshift=2mm},
  title={(a) 积分距离小的两个函数},samples=241,domain=0:1,
  legend style={at={(0.97,0.03)},anchor=south east}]
\addplot[draw=FigureBlue,line width=1.1pt] {sqrt(2)*sin(deg(0.5*pi*x))};
\addplot[draw=FigureYellow,line width=.9pt] {sqrt(2)*sin(deg(0.5*pi*x))+0.1*sqrt(2)*sin(deg(7.5*pi*x))};
\legend{$f=\psi_1$,$g=\psi_1+0.1\psi_8$}
\end{axis}
\begin{semilogyaxis}[at={(A.east)},anchor=west,xshift=11mm,xmin=0.4,xmax=10.6,ymin=0.5,ymax=1500,
  xtick={1,2,...,10},ytick={1,10,100,1000},xlabel={$j$},
  xlabel style={at={(axis description cs:1,0)},anchor=north west,yshift=2mm},
  title={(b) 坐标平方的权重},
  legend style={at={(0.97,0.14)},anchor=south east}]
\addplot[FigureBlue,line width=.9pt,mark=*,mark size=1.3pt] coordinates {(1,2.4674) (2,22.2066) (3,61.6850) (4,120.9027) (5,199.8595) (6,298.5555) (7,416.9908) (8,555.1652) (9,713.0789) (10,890.7318)};
\addplot[FigureGreen,line width=.9pt,dashed,mark=square*,mark size=1.2pt,mark options={solid}] coordinates {(1,1) (2,1) (3,1) (4,1) (5,1) (6,1) (7,1) (8,1) (9,1) (10,1)};
\legend{$H_K$：$1/\lambda_j$,$L^2$：$1$}
\end{semilogyaxis}
\end{tikzpicture}
